\documentclass[11pt]{article}

% -------------------------------------------------
% Packages
% -------------------------------------------------

\usepackage[margin=1in]{geometry}
\usepackage[T1]{fontenc}
\usepackage[utf8]{inputenc}
\usepackage{lmodern}
\usepackage{amsmath}
\usepackage{amssymb}
\usepackage{mathtools}
\usepackage{enumitem}

% -------------------------------------------------
% Formatting
% -------------------------------------------------

\setlength{\parindent}{0pt}
\setlength{\parskip}{0.45em}

\setlist[enumerate]{
    label=\alph*),
    leftmargin=2em,
    itemsep=1.2em
}

\pagestyle{plain}

\newcommand{\problem}[1]{%
    \vspace{1em}
    {\large\bfseries #1}
    \par
    \vspace{0.4em}
}

\newcommand{\E}{\mathbb{E}}
\newcommand{\Var}{\operatorname{Var}}
\newcommand{\KL}{D_{\mathrm{KL}}}

% -------------------------------------------------
% Document
% -------------------------------------------------

\begin{document}

\begin{center}
    {\LARGE \textbf{Problem Set 1}}\\[0.4em]
    {\large Physics 286/AP 286/MCB 198}\\[0.25em]
    {\large Fall 2026}\\[0.4em]
    {\large \textbf{Due September 23}}
\end{center}

\vspace{0.8em}

Please do not use LLMs or other AI tools for problems requiring
pencil-and-paper work. You are welcome to discuss the problems with your
classmates, but write your solutions independently after your discussions.
Please name the people with whom you discussed the problems in the solutions
you submit.

% =================================================
% Problem I.1
% =================================================

\problem{OPTIONAL FOR EVERYBODY: Problem I.1 The Hat Game and the Cost of Coordination}

Three players---Alice, Bob, and Charlie---are each randomly assigned either
a red or a blue hat. Each player can see everyone else's hat but not their
own. They are not allowed to communicate once the game begins, but they may
agree on a strategy beforehand.

After observing the other two hats, each player may either guess the color
of their own hat or pass. The team wins if at least one player guesses
correctly and no player guesses incorrectly.

\begin{enumerate}

\item Suppose each player guesses at random, independently, or suppose
that all players always pass. What is the probability that the team wins?

If all three guess randomly, each individual guess is correct with
probability \(1/2\). Since a single wrong guess makes the team lose, all
three have to be correct:

\[
P(\text{win})
=
\left(\frac12\right)^3
=
\frac18.
\]

If everyone passes, there is no correct guess, so

\[
P(\text{win})=0.
\]

\item Show that if the players coordinate in advance, they can use a
strategy that increases their probability of winning above chance.
What is the best strategy, and what is the resulting win probability?

Use the following rule:

\[
\begin{cases}
\text{if the other two hats match, guess the opposite color},\\
\text{if the other two hats differ, pass}.
\end{cases}
\]

For example, suppose the hats are \(RRB\). The player wearing blue sees
\(RR\), so they guess blue. The other two players each see one red and one
blue, so they pass. Thus the team wins.

The same thing happens whenever exactly one hat differs from the other two.
There are

\[
\binom31(2)=6
\]

such assignments. The only losing assignments are

\[
RRR,\qquad BBB.
\]

Therefore

\[
P(\text{win})=\frac68=\frac34.
\]

For optimality, think of the \(8\) hat assignments as the vertices of a cube.
If an assignment is winning, flip the hat of a player who guessed correctly.
The new assignment must be losing, since that player's observation does not
change but their guess becomes wrong.

Thus every winning vertex is adjacent to a losing vertex. Each losing vertex
has only \(3\) neighbors, so if \(W\) is the number of winning states and
\(L\) the number of losing states,

\[
W\leq 3L.
\]

Also

\[
W+L=8.
\]

Hence

\[
W\leq 3(8-W),
\]

so

\[
4W\leq 24
\]

and therefore

\[
W\leq 6.
\]

So the \(6/8=3/4\) strategy is optimal.

\item Let us quantify the information available to each player. How
many bits of uncertainty does each player have about their own hat?
How much information is available from seeing the other players' hats?

A player's own hat is a fair binary random variable, so

\[
H(\text{own})
=
-\frac12\log_2\frac12
-\frac12\log_2\frac12
=
1\text{ bit}.
\]

The other two hats have four equally likely possibilities,

\[
RR,\ RB,\ BR,\ BB,
\]

so their entropy is

\[
H(\text{other two})=\log_2 4=2\text{ bits}.
\]

But the hats are assigned independently. Therefore seeing the other two
does not tell a player anything about their own hat:

\[
H(\text{own}\mid\text{other two})
=
H(\text{own})
=
1.
\]

Thus

\[
I(\text{own};\text{other two})
=
H(\text{own})-H(\text{own}\mid\text{other two})
=
0.
\]

\item Design a variant of this game in which each player wears one of
four hat colors instead of two. Is there still a winning strategy that
performs better than chance? What principles from information theory
might be useful?

Label the four colors

\[
0,1,2,3\pmod 4.
\]

Beforehand, agree to act as though the total color sum should satisfy

\[
x_1+x_2+x_3=0\pmod 4.
\]

Each player sees the other two colors and guesses the unique color that
would make the sum \(0\pmod4\).

For example, player \(1\) guesses

\[
\hat x_1=-(x_2+x_3)\pmod4.
\]

All three guesses are correct exactly when the actual colors satisfy

\[
x_1+x_2+x_3=0\pmod4.
\]

For any choices of \(x_1,x_2\), exactly one of the four possible values of
\(x_3\) makes this true. Thus

\[
P(\text{win})=\frac14.
\]

By comparison, if all three independently guess one of four colors,

\[
P(\text{all correct})
=
\left(\frac14\right)^3
=
\frac1{64}.
\]

So coordination again helps substantially. The useful ideas are
\(q\)-ary codes, modular checksums, parity constraints, and Hamming distance.

\item Suppose the referee tells the players the parity---that is,
whether the number of red hats is even or odd. How does this change
the information content of each observation, and how would you modify
the strategy?

Without the parity bit, after seeing the other two hats a player still has

\[
H(\text{own}\mid\text{other two})=1\text{ bit}
\]

of uncertainty.

Now suppose the parity is also given. If a player sees the other two hats,
they know whether those two contain an even or odd number of red hats.
Comparing that with the announced total parity determines their own hat.

For example, if the total parity is odd and the player sees two red hats,
then the visible hats contribute an even number of reds, so their own hat
must be red.

Therefore

\[
H(\text{own}\mid\text{other two, parity})=0.
\]

The parity supplies exactly the missing information needed to determine
the player's hat. All three players can infer their hats exactly, so

\[
P(\text{win})=1.
\]

\end{enumerate}

% =================================================
% Problem I.2
% =================================================

\problem{Problem I.2: Evaluating Kullback--Leibler Distances Between Gaussian Distributions}

Consider

\[
p(x)=\frac{1}{\sqrt{2\pi\sigma^2}}
\exp\left[-\frac{(x-\mu_1)^2}{2\sigma^2}\right]
\]

and

\[
q(x)=\frac{1}{\sqrt{2\pi}}
\exp\left[-\frac{(x-\mu_2)^2}{2}\right].
\]

\begin{enumerate}

\item Evaluate the Kullback--Leibler distance between \(p(x)\) and \(q(x)\).

Start with

\[
\KL(p\|q)
=
\int p(x)\log\frac{p(x)}{q(x)}\,dx.
\]

The log ratio is

\[
\log\frac{p(x)}{q(x)}
=
-\frac12\log\sigma^2
-\frac{(x-\mu_1)^2}{2\sigma^2}
+\frac{(x-\mu_2)^2}{2}.
\]

Take expectation under \(p\):

\[
\KL(p\|q)
=
-\frac12\log\sigma^2
-\frac{1}{2\sigma^2}
\E_p[(X-\mu_1)^2]
+\frac12\E_p[(X-\mu_2)^2].
\]

Since \(X\sim N(\mu_1,\sigma^2)\),

\[
\E_p[(X-\mu_1)^2]=\sigma^2.
\]

Also,

\[
X-\mu_2
=
(X-\mu_1)+(\mu_1-\mu_2),
\]

so

\[
\E_p[(X-\mu_2)^2]
=
\sigma^2+(\mu_1-\mu_2)^2.
\]

Substituting,

\begin{align*}
\KL(p\|q)
&=
-\frac12\log\sigma^2
-\frac12
+\frac12\left[
\sigma^2+(\mu_1-\mu_2)^2
\right] \\
&=
\frac12
\left[
\sigma^2+(\mu_1-\mu_2)^2-1-\log\sigma^2
\right].
\end{align*}

\item Find the value of \(\mu_1\) that minimizes this distance for fixed
\(\mu_2\) and \(\sigma\).

Only the term

\[
\frac12(\mu_1-\mu_2)^2
\]

depends on \(\mu_1\). Differentiate:

\[
\frac{\partial \KL}{\partial\mu_1}
=
\mu_1-\mu_2.
\]

Set this equal to zero:

\[
\mu_1-\mu_2=0,
\]

so

\[
\mu_1=\mu_2.
\]

The second derivative is

\[
\frac{\partial^2\KL}{\partial\mu_1^2}=1>0,
\]

so this is indeed the minimum.

\end{enumerate}

% =================================================
% Problem I.3
% =================================================

\problem{Problem I.3: Mo's Variational Compression Scheme}

Jen sends data to Mo. Each data point \(x\in\mathbb N\) is independently drawn
from a Poisson distribution with mean \(\lambda>0\):

\[
P_\lambda(x)=\frac{\lambda^x e^{-\lambda}}{x!}.
\]

Mo assumes instead a Gaussian model

\[
Q_{\mu,\sigma}(x)=\frac{1}{\sqrt{2\pi\sigma^2}}
\exp\left[-\frac{(x-\mu)^2}{2\sigma^2}\right]
\]

and minimizes

\[
\KL(P_\lambda\|Q_{\mu,\sigma})
=
\sum_{x=0}^{\infty}
P_\lambda(x)
\log\frac{P_\lambda(x)}{Q_{\mu,\sigma}(x)}.
\]

\begin{enumerate}

\item Explain why minimizing this KL divergence is sensible from the
perspective of data compression.

If Mo uses \(Q\) to encode data actually generated from \(P\), the expected
code length is proportional to

\[
\E_P[-\log Q(X)].
\]

The optimal expected length, if Mo knew the true distribution, would be

\[
\E_P[-\log P(X)].
\]

The difference is

\begin{align*}
\E_P[-\log Q(X)]
-
\E_P[-\log P(X)]
&=
\E_P\left[
\log\frac{P(X)}{Q(X)}
\right] \\
&=
\KL(P\|Q).
\end{align*}

Thus KL divergence is exactly the extra expected coding cost from using
\(Q\) instead of \(P\). Minimizing it minimizes the compression overhead.

\item Show that the KL divergence, up to a constant independent of
\(\mu,\sigma\), has the stated form.

Write

\[
\KL(P_\lambda\|Q_{\mu,\sigma})
=
\E_{P_\lambda}
\left[
\log P_\lambda(X)
-
\log Q_{\mu,\sigma}(X)
\right].
\]

The first term does not depend on \(\mu\) or \(\sigma\), so it can be absorbed
into a constant.

For the Gaussian,

\[
-\log Q_{\mu,\sigma}(X)
=
\frac12\log(2\pi\sigma^2)
+
\frac{(X-\mu)^2}{2\sigma^2}.
\]

Hence

\[
\KL(P_\lambda\|Q_{\mu,\sigma})
=
\E_{P_\lambda}
\left[
\frac{(X-\mu)^2}{2\sigma^2}
+
\frac12\log(2\pi\sigma^2)
\right]
+\text{const.}
\]

\item Using \(\E[X]=\lambda\) and \(\Var(X)=\lambda\), derive the optimal
\(\mu\) and \(\sigma\).

First expand

\[
X-\mu
=
(X-\lambda)+(\lambda-\mu).
\]

Therefore

\begin{align*}
\E[(X-\mu)^2]
&=
\E[(X-\lambda)^2]
+2(\lambda-\mu)\E[X-\lambda]
+(\lambda-\mu)^2 \\
&=
\lambda+(\lambda-\mu)^2,
\end{align*}

because

\[
\E[X-\lambda]=0,
\qquad
\E[(X-\lambda)^2]=\Var(X)=\lambda.
\]

Ignoring constants, the objective is therefore

\[
F(\mu,\sigma)
=
\frac{\lambda+(\lambda-\mu)^2}{2\sigma^2}
+
\frac12\log(2\pi\sigma^2).
\]

Differentiate with respect to \(\mu\):

\[
\frac{\partial F}{\partial\mu}
=
\frac{\mu-\lambda}{\sigma^2}.
\]

Thus

\[
\mu=\lambda.
\]

Now substitute this back in:

\[
F(\lambda,\sigma)
=
\frac{\lambda}{2\sigma^2}
+
\frac12\log(2\pi\sigma^2).
\]

It is a little cleaner to let

\[
v=\sigma^2.
\]

Then

\[
F(v)=\frac{\lambda}{2v}+\frac12\log(2\pi v).
\]

Differentiate:

\[
\frac{dF}{dv}
=
-\frac{\lambda}{2v^2}
+
\frac{1}{2v}.
\]

Set equal to zero:

\[
-\frac{\lambda}{2v^2}
+
\frac{1}{2v}
=0.
\]

Multiply by \(2v^2\):

\[
-\lambda+v=0.
\]

Thus

\[
v=\lambda,
\]

so

\[
\sigma^2=\lambda
\]

and therefore

\[
\sigma=\sqrt{\lambda}.
\]

So the optimal parameters are

\[
\mu^*=\lambda,
\qquad
(\sigma^*)^2=\lambda.
\]

\item Briefly interpret your result.

The minimizing Gaussian has

\[
\text{mean}=\lambda
\]

and

\[
\text{variance}=\lambda,
\]

which are exactly the mean and variance of the Poisson distribution.

So in this KL direction, the best Gaussian in the allowed family matches
the first two moments of the Poisson. For large \(\lambda\), this is also
the usual Gaussian approximation to a Poisson random variable.

\end{enumerate}

% =================================================
% Problem I.4
% =================================================

\problem{Problem I.4 Running a Casino}

Suppose a coin has \(P(H)=p\) and \(P(T)=1-p\). A bet of \(x\) dollars on
heads pays \(g_Hx\) if heads occurs, and a bet of \(x\) dollars on tails pays
\(g_Tx\) if tails occurs. Define fair odds that do not bankrupt the casino in
the long run.

For a fair bet, the expected payout should equal the amount wagered.

For a heads bet,

\[
\E[\text{payout}]
=
p(g_Hx)+(1-p)(0).
\]

Set this equal to the original stake \(x\):

\[
p g_H x=x.
\]

Assuming \(x\neq0\),

\[
p g_H=1,
\]

so

\[
g_H=\frac1p.
\]

Similarly, for tails,

\[
(1-p)g_Tx=x,
\]

which gives

\[
g_T=\frac{1}{1-p}.
\]

Thus fair payout multipliers are

\[
g_H=\frac1p,
\qquad
g_T=\frac{1}{1-p}.
\]

With these odds the expected net gain is zero for both the gambler and
casino.

% =================================================
% Problem I.5
% =================================================

\problem{Problem I.5 Solving a Poorly Posed Problem Using Maximum Entropy}

\textbf{Inferring models of cars.}

Two models of cars of lengths \(\ell_1\) and \(\ell_2\) are loaded end-to-end
onto a barge. Let \(n_1,n_2\) be the unknown counts, and suppose the only
aggregate information is the ship length \(L\).

Let

\[
P_{n_1,n_2}=P(n_1,n_2)
\]

be a probability distribution on nonnegative integer pairs
\((n_1,n_2)\).

The constraints are normalization,

\[
\sum_{n_1,n_2}P_{n_1,n_2}=1,
\]

and the known mean total length,

\[
\sum_{n_1,n_2}
P_{n_1,n_2}
(\ell_1n_1+\ell_2n_2)
=
L.
\]

We maximize the entropy

\[
S
=
-\sum_{n_1,n_2}
P_{n_1,n_2}\log P_{n_1,n_2}.
\]

Introduce Lagrange multipliers \(\alpha,\beta\):

\[
\mathcal L
=
-\sum P\log P
-\alpha\left(\sum P-1\right)
-\beta
\left[
\sum P(\ell_1n_1+\ell_2n_2)-L
\right].
\]

Differentiate with respect to one probability:

\[
\frac{\partial\mathcal L}{\partial P_{n_1,n_2}}
=
-\log P_{n_1,n_2}
-1
-\alpha
-\beta(\ell_1n_1+\ell_2n_2)
=0.
\]

Thus

\[
\log P_{n_1,n_2}
=
-1-\alpha
-\beta(\ell_1n_1+\ell_2n_2),
\]

or, after absorbing the extra constant into the normalization,

\[
P(n_1,n_2)
=
\frac{
e^{-\beta(\ell_1n_1+\ell_2n_2)}
}{Z}.
\]

The partition function is

\begin{align*}
Z
&=
\sum_{n_1=0}^{\infty}
\sum_{n_2=0}^{\infty}
e^{-\beta\ell_1n_1}
e^{-\beta\ell_2n_2} \\
&=
\left(
\sum_{n_1=0}^{\infty}e^{-\beta\ell_1n_1}
\right)
\left(
\sum_{n_2=0}^{\infty}e^{-\beta\ell_2n_2}
\right).
\end{align*}

Each sum is geometric:

\[
\sum_{n=0}^{\infty}r^n=\frac{1}{1-r}.
\]

Hence

\[
Z
=
\frac{1}{
(1-e^{-\beta\ell_1})
(1-e^{-\beta\ell_2})
}.
\]

The mean of each count is the mean of a geometric distribution:

\[
\E[n_i]
=
\frac{e^{-\beta\ell_i}}
{1-e^{-\beta\ell_i}}
=
\frac{1}{e^{\beta\ell_i}-1}.
\]

Therefore the mean length constraint becomes

\[
\ell_1\E[n_1]+\ell_2\E[n_2]=L,
\]

so

\[
\frac{\ell_1}{e^{\beta\ell_1}-1}
+
\frac{\ell_2}{e^{\beta\ell_2}-1}
=
L.
\]

This equation determines \(\beta\).

Since

\[
P(n_1,n_2)
\propto
e^{-\beta(\ell_1n_1+\ell_2n_2)},
\]

the probability decreases as either \(n_1\) or \(n_2\) increases. Therefore,
with only the average-length constraint, the literal mode is

\[
(n_1,n_2)=(0,0).
\]

This is a sign that the inference problem is poorly posed: imposing only an
average total length is not the same as saying that every barge is filled to
length \(L\).

For

\[
(\ell_1,\ell_2,L)=(3,11,2000),
\]

solving

\[
\frac{3}{e^{3\beta}-1}
+
\frac{11}{e^{11\beta}-1}
=
2000
\]

gives approximately

\[
\beta
\approx
9.96518\times10^{-4}.
\]

Then

\[
\E[n_1]
=
\frac{1}{e^{3\beta}-1}
\approx333.998
\]

and

\[
\E[n_2]
=
\frac{1}{e^{11\beta}-1}
\approx90.728.
\]

If instead the barge must be exactly filled in every realization, then the
support should be restricted to pairs satisfying

\[
\ell_1n_1+\ell_2n_2=L.
\]

That would be a different maximum-entropy problem.

% =================================================
% Problem I.6
% =================================================

\problem{Problem I.6 Statistical Mechanics of Binding and Unbinding of Transcription Factors}

Consider a particle with states \(i=1,\ldots,N\), probabilities \(p_i\), and
energies \(E_i\), with fixed average energy
\(\sum_i p_iE_i=\langle E\rangle\).

\begin{enumerate}

\item What distribution maximizes entropy? Define the partition function
\(Z\), and evaluate \(-\partial\ln Z/\partial\beta\).

We maximize

\[
S=-\sum_i p_i\log p_i
\]

subject to

\[
\sum_i p_i=1
\]

and

\[
\sum_i p_iE_i=\langle E\rangle.
\]

Using multipliers \(\alpha,\beta\),

\[
\mathcal L
=
-\sum_i p_i\log p_i
-\alpha\left(\sum_i p_i-1\right)
-\beta\left(\sum_i p_iE_i-\langle E\rangle\right).
\]

Differentiate:

\[
\frac{\partial\mathcal L}{\partial p_i}
=
-\log p_i-1-\alpha-\beta E_i=0.
\]

Therefore

\[
p_i\propto e^{-\beta E_i}.
\]

Normalizing gives

\[
p_i=\frac{e^{-\beta E_i}}{Z},
\]

where

\[
Z=\sum_i e^{-\beta E_i}.
\]

Now

\[
\frac{\partial\ln Z}{\partial\beta}
=
\frac1Z\frac{\partial Z}{\partial\beta}.
\]

Also

\[
\frac{\partial Z}{\partial\beta}
=
-\sum_i E_i e^{-\beta E_i}.
\]

Hence

\begin{align*}
-\frac{\partial\ln Z}{\partial\beta}
&=
\frac{1}{Z}
\sum_i E_i e^{-\beta E_i} \\
&=
\sum_i E_i
\frac{e^{-\beta E_i}}{Z} \\
&=
\sum_i p_iE_i \\
&=
\langle E\rangle.
\end{align*}

\item A cell of volume \(V\) contains transcription factor at concentration
\(T\). The gene is unbound or bound by one molecule. A free molecule has
energy \(E_0\); a bound molecule has energy \(E_1\). Divide the cell into
boxes of volume \(1/C\), with \(CV\gg TV\), and derive
\(p_{\mathrm{bound}}\).

Let

\[
M=CV
\]

be the number of boxes and

\[
N=TV
\]

the number of transcription factor molecules.

If the gene is unbound, all \(N\) molecules are free. The number of ways to
place them is

\[
\binom{M}{N},
\]

so the statistical weight is

\[
w_0
=
\binom{M}{N}
e^{-\beta N E_0}.
\]

If the gene is bound, one molecule is on the gene and \(N-1\) are free:

\[
w_1
=
\binom{M}{N-1}
e^{-\beta[(N-1)E_0+E_1]}.
\]

Take the ratio:

\begin{align*}
\frac{w_1}{w_0}
&=
\frac{\binom{M}{N-1}}{\binom{M}{N}}
e^{-\beta(E_1-E_0)} \\
&=
\frac{N}{M-N+1}
e^{-\beta(E_1-E_0)}.
\end{align*}

Since \(M\gg N\),

\[
M-N+1\simeq M.
\]

Thus

\[
\frac{w_1}{w_0}
\simeq
\frac{N}{M}
e^{-\beta(E_1-E_0)}.
\]

Since

\[
\frac{N}{M}
=
\frac{TV}{CV}
=
\frac{T}{C},
\]

we get

\[
\frac{w_1}{w_0}
\simeq
\frac{T}{C}
e^{-\beta(E_1-E_0)}.
\]

Let

\[
r=
\frac{T}{C}
e^{-\beta(E_1-E_0)}.
\]

Then

\[
p_{\mathrm{bound}}
=
\frac{w_1}{w_0+w_1}
=
\frac{r}{1+r}.
\]

Define

\[
K
=
C e^{\beta(E_1-E_0)}.
\]

Since \(r=T/K\),

\[
p_{\mathrm{bound}}
=
\frac{T}{T+K}.
\]

\item Suppose the gene can bind one or two molecules. What is the average
number bound as a function of \(T\)?

Define

\[
z
=
\frac{T}{C}
e^{-\beta(E_1-E_0)}.
\]

The relative statistical weights for having \(0,1,2\) bound molecules are

\[
1,\qquad z,\qquad z^2.
\]

Therefore

\[
Z_{\text{bind}}
=
1+z+z^2.
\]

The average number bound is

\[
\langle n_{\mathrm{bound}}\rangle
=
\frac{
0(1)+1(z)+2(z^2)
}{
1+z+z^2
}.
\]

Thus

\[
\langle n_{\mathrm{bound}}\rangle
=
\frac{z+2z^2}{1+z+z^2}.
\]

\item Suppose the second molecule binds with energy \(E_2\neq E_1\). What
is the average number bound now?

Let

\[
z_1
=
\frac{T}{C}
e^{-\beta(E_1-E_0)}
\]

for binding the first molecule and

\[
z_2
=
\frac{T}{C}
e^{-\beta(E_2-E_0)}
\]

for binding the second.

The relative weights are now

\[
1,
\qquad
z_1,
\qquad
z_1z_2.
\]

Hence

\[
Z_{\text{bind}}
=
1+z_1+z_1z_2.
\]

Then

\[
\langle n_{\mathrm{bound}}\rangle
=
\frac{
0(1)+1(z_1)+2(z_1z_2)
}{
1+z_1+z_1z_2
}.
\]

So

\[
\langle n_{\mathrm{bound}}\rangle
=
\frac{
z_1+2z_1z_2
}{
1+z_1+z_1z_2
}.
\]

\end{enumerate}

% =================================================
% Problem I.7
% =================================================

\problem{Problem I.7: Finding Key Amino Acids in Protein Families}

The files \texttt{sequencefamily*.txt} contain 100 sequences for each of four
families A--D, and \texttt{sequencefreqs*.txt} contain amino-acid frequencies
at the 10 variable sites.

\begin{enumerate}

\item For each site and family, calculate \(H(a_i\mid f)\), and also
calculate \(H(a_i)\) across all families.

For a particular family \(f\), let

\[
p_f(a;i)
\]

be the frequency of amino acid \(a\) at site \(i\).

Then the entropy at that site within family \(f\) is

\[
H(a_i\mid f)
=
-\sum_a p_f(a;i)\log_2 p_f(a;i).
\]

Since the four families have equal size, the overall amino-acid frequency at
site \(i\) is

\[
p(a;i)
=
\frac14
\sum_f p_f(a;i).
\]

Then

\[
H(a_i)
=
-\sum_a p(a;i)\log_2p(a;i).
\]

So for each of the \(10\) sites, the calculation is:

\[
p_f(a;i)
\longrightarrow
H(a_i\mid f)
\]

for each family, followed by

\[
p(a;i)
=
\frac14\sum_f p_f(a;i)
\longrightarrow
H(a_i).
\]

The numerical values require the sequence/frequency files, which were not
included with this problem set.

\item Compute the mutual information \(I(a_i;F)\). What are its possible
minimum and maximum values?

The conditional entropy given the family is

\[
H(a_i\mid F)
=
\sum_f P(f)H(a_i\mid f).
\]

Since all four families are equally likely,

\[
P(f)=\frac14,
\]

so

\[
H(a_i\mid F)
=
\frac14\sum_f H(a_i\mid f).
\]

The mutual information is

\[
I(a_i;F)
=
H(a_i)-H(a_i\mid F).
\]

The minimum is

\[
I(a_i;F)=0,
\]

which happens when knowing the family gives no information about the residue.

The maximum cannot exceed the entropy of the family label:

\[
I(a_i;F)\leq H(F).
\]

Since there are four equally likely families,

\[
H(F)
=
-\sum_{f=1}^4
\frac14\log_2\frac14
=
2\text{ bits}.
\]

Therefore

\[
0\leq I(a_i;F)\leq2\text{ bits}.
\]

Numerical values for each site require the supplied data files.

\item A new protein has phenylalanine at site 10. To which family should
it be assigned, and with what probability?

We want

\[
P(f\mid a_{10}=\mathrm{Phe}).
\]

By Bayes' rule,

\[
P(f\mid\mathrm{Phe})
=
\frac{
P(\mathrm{Phe}\mid f)P(f)
}{
P(\mathrm{Phe})
}.
\]

Since the priors are equal,

\[
P(f)=\frac14.
\]

Also,

\[
P(\mathrm{Phe})
=
\sum_g
P(\mathrm{Phe}\mid g)P(g)
=
\frac14
\sum_g
P(\mathrm{Phe}\mid g).
\]

The factors of \(1/4\) cancel:

\[
P(f\mid a_{10}=\mathrm{Phe})
=
\frac{
P(a_{10}=\mathrm{Phe}\mid f)
}{
\sum_g
P(a_{10}=\mathrm{Phe}\mid g)
}.
\]

So the protein should be assigned to whichever family has the largest
site-10 phenylalanine frequency.

The specific family and probability cannot be determined without the
site-10 frequency data.

\item Divide the sites into conserved, specificity-determining, and
neither categories.

A conserved site should have little uncertainty overall, so

\[
H(a_i)\approx0.
\]

A specificity-determining site should vary across families but be predictable
once the family is known. Equivalently,

\[
I(a_i;F)
\]

should be large.

Thus the rough classification is

\[
\begin{array}{ll}
\text{conserved:}
&
H(a_i)\text{ small},
\\[0.4em]
\text{specificity-determining:}
&
I(a_i;F)\text{ large},
\\[0.4em]
\text{neither:}
&
H(a_i)\text{ not small and }I(a_i;F)\text{ small}.
\end{array}
\]

Exact site numbers require the missing sequence/frequency files.

\end{enumerate}

% =================================================
% Problem I.8
% =================================================

\problem{Problem I.8 Determining the Capacity of a Noisy Communication Channel}

The joint distribution of input \(X\in\{a,b,c,d\}\) and output
\(Y\in\{a,b,c,d\}\) is

\[
\begin{array}{c|cccc}
    & x=a & x=b & x=c & x=d \\
    \hline
    y=a & \frac18  & \frac1{16} & \frac1{16} & \frac14 \\
    y=b & \frac1{16} & \frac18  & \frac1{16} & 0 \\
    y=c & \frac1{32} & \frac1{32} & \frac1{16} & 0 \\
    y=d & \frac1{32} & \frac1{32} & \frac1{16} & 0
\end{array}
\]

\begin{enumerate}

\item Find \(P_X\) and \(H(X)\).

Sum down each column.

For \(x=a\),

\[
P_X(a)
=
\frac18+\frac1{16}+\frac1{32}+\frac1{32}
=
\frac14.
\]

Similarly,

\[
P_X(b)=\frac14,
\qquad
P_X(c)=\frac14,
\qquad
P_X(d)=\frac14.
\]

Thus

\[
P_X
=
\left(
\frac14,\frac14,\frac14,\frac14
\right).
\]

Then

\begin{align*}
H(X)
&=
-\sum_xP_X(x)\log_2P_X(x) \\
&=
-4\left(\frac14\right)\log_2\frac14 \\
&=
2\text{ bits}.
\end{align*}

\item Find \(P_Y\) and \(H(Y)\).

Sum across each row.

For \(y=a\),

\[
P_Y(a)
=
\frac18+\frac1{16}+\frac1{16}+\frac14
=
\frac12.
\]

For \(y=b\),

\[
P_Y(b)
=
\frac1{16}+\frac18+\frac1{16}
=
\frac14.
\]

For \(y=c\),

\[
P_Y(c)
=
\frac1{32}+\frac1{32}+\frac1{16}
=
\frac18.
\]

Similarly,

\[
P_Y(d)=\frac18.
\]

Therefore

\[
P_Y
=
\left(
\frac12,\frac14,\frac18,\frac18
\right).
\]

Now

\begin{align*}
H(Y)
&=
-\frac12\log_2\frac12
-\frac14\log_2\frac14
-2\left(\frac18\log_2\frac18\right) \\
&=
\frac12+\frac12+\frac34 \\
&=
\frac74.
\end{align*}

Thus

\[
H(Y)=1.75\text{ bits}.
\]

\item Compute \(H(X,Y)\).

The nonzero probabilities in the table are:

\[
1\times\frac14,
\qquad
2\times\frac18,
\qquad
6\times\frac1{16},
\qquad
4\times\frac1{32}.
\]

Therefore

\begin{align*}
H(X,Y)
&=
-\frac14\log_2\frac14
-2\left(\frac18\log_2\frac18\right) \\
&\qquad
-6\left(\frac1{16}\log_2\frac1{16}\right)
-4\left(\frac1{32}\log_2\frac1{32}\right).
\end{align*}

Compute each piece:

\[
-\frac14\log_2\frac14
=
\frac12,
\]

\[
-2\left(\frac18\log_2\frac18\right)
=
2\left(\frac38\right)
=
\frac34,
\]

\[
-6\left(\frac1{16}\log_2\frac1{16}\right)
=
6\left(\frac4{16}\right)
=
\frac32,
\]

and

\[
-4\left(\frac1{32}\log_2\frac1{32}\right)
=
4\left(\frac5{32}\right)
=
\frac58.
\]

Hence

\[
H(X,Y)
=
\frac12+\frac34+\frac32+\frac58
=
\frac{27}{8}
=
3.375\text{ bits}.
\]

\item Compute \(H(Y\mid X)\).

Use

\[
H(Y\mid X)
=
H(X,Y)-H(X).
\]

Therefore

\[
H(Y\mid X)
=
\frac{27}{8}-2
=
\frac{27}{8}-\frac{16}{8}
=
\frac{11}{8}.
\]

Thus

\[
H(Y\mid X)=1.375\text{ bits}.
\]

\item Compute \(I(X;Y)\).

Use

\[
I(X;Y)
=
H(X)+H(Y)-H(X,Y).
\]

Substitute:

\begin{align*}
I(X;Y)
&=
2+\frac74-\frac{27}{8} \\
&=
\frac{16}{8}+\frac{14}{8}-\frac{27}{8} \\
&=
\frac38.
\end{align*}

So

\[
I(X;Y)=0.375\text{ bits}.
\]

\item Give a lower-bound estimate of the channel capacity \(C\).

Channel capacity is

\[
C
=
\max_{P_X} I(X;Y).
\]

The supplied distribution uses the uniform input distribution and gives

\[
I(X;Y)=\frac38.
\]

Since the maximum over all input distributions must be at least as large as
the value for this particular one,

\[
C\geq\frac38\text{ bits/use}.
\]

\end{enumerate}

% =================================================
% Problem I.9
% =================================================

\problem{Problem I.9 Optimizing the Source Code to Match the Capacity of an Information Channel}

A source produces A, B, and C with probabilities \(P,Q,Q\), where

\[
P+2Q=1.
\]

A is transmitted perfectly. B and C are transmitted correctly with
probability \(p\) and interchanged with probability \(q\), so \(p+q=1\).

\begin{enumerate}

\item Compute \(H(X)\) and \(H(X\mid Y)\).

The input distribution is

\[
(P,Q,Q).
\]

Thus

\begin{align*}
H(X)
&=
-P\log_2P
-Q\log_2Q
-Q\log_2Q \\
&=
-P\log_2P
-2Q\log_2Q.
\end{align*}

Now define the binary entropy associated with confusion between B and C:

\[
h
=
-p\log_2p-q\log_2q.
\]

If the output is A, then the input must have been A, so

\[
H(X\mid Y=A)=0.
\]

If the output is B or C, there remains a binary uncertainty between B and C
with entropy \(h\).

By symmetry the output probabilities are also

\[
P_Y=(P,Q,Q).
\]

Therefore

\begin{align*}
H(X\mid Y)
&=
P(Y=A)(0)
+
P(Y=B)h
+
P(Y=C)h \\
&=
Qh+Qh \\
&=
2Qh.
\end{align*}

So

\[
H(X)
=
-P\log_2P-2Q\log_2Q
\]

and

\[
H(X\mid Y)=2Qh.
\]

\item Compute and maximize the transmission rate subject to \(P+2Q=1\).

The transmission rate is the mutual information

\[
I(X;Y)
=
H(X)-H(X\mid Y).
\]

Thus

\[
I(X;Y)
=
-P\log_2P
-2Q\log_2Q
-2Qh.
\]

We maximize this subject to

\[
P+2Q=1.
\]

Introduce a multiplier \(\lambda\):

\[
\mathcal L
=
-P\log_2P
-2Q\log_2Q
-2Qh
+
\lambda(P+2Q-1).
\]

Differentiate with respect to \(P\):

\[
\frac{\partial\mathcal L}{\partial P}
=
-\log_2P-\frac1{\ln2}+\lambda
=
0.
\]

Therefore

\[
\lambda
=
\log_2P+\frac1{\ln2}.
\]

Differentiate with respect to \(Q\):

\[
\frac{\partial\mathcal L}{\partial Q}
=
-2\log_2Q
-\frac2{\ln2}
-2h
+2\lambda
=
0.
\]

Divide by \(2\):

\[
-\log_2Q
-\frac1{\ln2}
-h
+\lambda
=
0.
\]

Thus

\[
\lambda
=
\log_2Q+\frac1{\ln2}+h.
\]

Equating the two expressions for \(\lambda\),

\[
\log_2P
=
\log_2Q+h.
\]

Therefore

\[
\log_2\frac{P}{Q}=h,
\]

so

\[
\frac{P}{Q}=2^h.
\]

Hence

\[
P=2^hQ.
\]

Use the normalization condition:

\[
2^hQ+2Q=1.
\]

Thus

\[
Q
=
\frac{1}{2^h+2}.
\]

Equivalently,

\[
Q
=
\frac{2^{-h}}{1+2^{1-h}}.
\]

Then

\[
P
=
2^hQ
=
\frac{2^h}{2^h+2}
=
\frac{1}{1+2^{1-h}}.
\]

So

\[
P^*
=
\frac{1}{1+2^{1-h}},
\qquad
Q^*
=
\frac{2^{-h}}{1+2^{1-h}}.
\]

\item Using the optimum, compute the channel capacity.

Start with

\[
I
=
-P\log_2P
-2Q\log_2Q
-2Qh.
\]

At the optimum,

\[
\frac{P}{Q}=2^h,
\]

so

\[
\log_2P
=
\log_2Q+h.
\]

Substitute this into \(I\):

\begin{align*}
I
&=
-P(\log_2Q+h)
-2Q\log_2Q
-2Qh \\
&=
-(P+2Q)\log_2Q
-(P+2Q)h.
\end{align*}

Since

\[
P+2Q=1,
\]

this becomes

\[
I
=
-\log_2Q-h.
\]

Using

\[
Q
=
\frac{2^{-h}}{1+2^{1-h}},
\]

we get

\begin{align*}
-\log_2Q
&=
-\log_2
\left(
\frac{2^{-h}}{1+2^{1-h}}
\right) \\
&=
h+\log_2(1+2^{1-h}).
\end{align*}

Therefore

\[
C
=
\log_2(1+2^{1-h}),
\]

where

\[
h=-p\log_2p-q\log_2q.
\]

\item Check \(p=1\) and \(p=1/2\).

First let

\[
p=1,
\qquad
q=0.
\]

Then

\[
h=0.
\]

So

\[
P^*
=
\frac{1}{1+2}
=
\frac13
\]

and

\[
Q^*
=
\frac{1}{1+2}
=
\frac13.
\]

Thus the three symbols are used equally often, as expected for a noiseless
three-symbol channel.

The capacity is

\[
C
=
\log_2(1+2)
=
\log_2 3.
\]

Now let

\[
p=q=\frac12.
\]

Then

\[
h
=
-\frac12\log_2\frac12
-\frac12\log_2\frac12
=
1.
\]

Therefore

\[
P^*
=
\frac{1}{1+2^{0}}
=
\frac12
\]

and

\[
Q^*
=
\frac{2^{-1}}{1+2^0}
=
\frac14.
\]

The capacity becomes

\[
C
=
\log_2(1+2^0)
=
\log_2 2
=
1\text{ bit/use}.
\]

This makes sense because when \(p=q=1/2\), the channel cannot distinguish B
from C. It can only reliably distinguish

\[
A
\]

from

\[
\text{``B or C''},
\]

which leaves an effectively binary channel.

\end{enumerate}

\end{document}